\documentclass[border=10pt]{standalone}
\usepackage[american voltages, american currents]{circuitikz}
\usepackage{amsmath}
\renewcommand{\familydefault}{\sfdefault}

\begin{document}
\begin{circuitikz}[line width=0.8pt, every node/.style={font=\sffamily}]

% T-equivalent circuit for opposing case
% La = L1 + M (series, left branch)
\draw (0,2) to[L, l=$L_a{=}L_1{+}M$, i>^=$i_1$] (3,2);

% Lc = -M (shunt element, center)
\draw (3,2) to[L, l_=$L_c{=}{-}M$] (3,0);

% Lb = L2 + M (series, right branch)
\draw (3,2) to[L, l=$L_b{=}L_2{+}M$, i>^=$i_2$] (6,2);

% Ground/reference connections
\draw (0,0) -- (6,0);
\draw (0,0) -- (0,0.2);
\draw (6,0) -- (6,0.2);

% Port labels
\node[above] at (0, 2.2) {Port 1};
\node[above] at (6, 2.2) {Port 2};

% Junction dot
\filldraw (3,2) circle (2pt);

\end{circuitikz}
\end{document}
